\documentclass[11pt]{article}
\pdfmapfile{=lm.map}
\pdfmapfile{+cm.map}
\pdfmapfile{+symbols.map}
\usepackage[T1]{fontenc}
\usepackage{lmodern,amsmath,amssymb,amsthm,booktabs,array,geometry,microtype,xcolor,hyperref}
\geometry{margin=1in}
\hypersetup{pdftitle={Two Endpoint Counterexamples to the Printed Erdos-Deep Triple Classification},pdfauthor={GPT-6 Astra},colorlinks=true,linkcolor=blue!45!black,urlcolor=blue!45!black,citecolor=blue!45!black}
\newtheorem{theorem}{Theorem}
\newtheorem{lemma}[theorem]{Lemma}
\newcommand{\Z}{\mathbb Z}
\newcommand{\AP}{\operatorname{AP}}
\newcommand{\dist}{\operatorname{dist}}
\newcommand{\ord}{\operatorname{ord}}
\title{Two Endpoint Counterexamples to the Printed\\Erd\H{o}s-Deep Triple Classification}
\author{GPT-6 Astra\\\small Project AI author label}
\date{14 October 2026}
\begin{document}
\maketitle
\begin{abstract}
We give primitive Erd\H{o}s-deep families of three modular arithmetic
progressions with length triples $(10,9,5)$ and $(20,5,5)$, respectively in
$\Z_{34}$ and $\Z_{48}$. Each constituent is itself Erd\H{o}s-deep.
Neither length triple occurs in the list in Dukes--Gaede's 2023 Conjecture~1.
Both constructions use a constituent of length
$\lfloor\ord(g)/2\rfloor+1$, permitted by the standing convention in
Gaede's thesis, Section~2.2. They do not refute the more restrictive variant
requiring every length to be at most $\lfloor\ord(g)/2\rfloor$.
Explicit pair counts and Lean~4 kernel-checked certificates establish the
examples. The source-convention discrepancy and the limits of the conclusion
are recorded precisely. No absolute priority or complete corrected
classification is claimed.
\end{abstract}

\section{Definitions and the endpoint}
For $n\ge1$, set
\[
 \AP_n(g,k)=\{jg\bmod n:0\le j<k\},\qquad
 \dist_n(x,y)=\min((x-y)\bmod n,(y-x)\bmod n).
\]
For a family $\mathcal F=(A_1,A_2,A_3)$, let $m_{\mathcal F}(d)$ count
unordered pairs at distance $d$, summed over the three constituents.
No pairs between different constituents are counted. The family is
\emph{Erd\H{o}s-deep} with parameter $K$ if its positive distance
multiplicities are precisely $1,\ldots,K-1$, each once. It is
\emph{primitive} if $\gcd(n,g_1,g_2,g_3)=1$.

Write $q=\ord(g)=n/\gcd(n,g)$. The constituent condition used here, with $k\ge3$, is
\begin{equation}\label{eq:half}
 k\le \lfloor q/2\rfloor+1,
 \quad\text{equivalently}\quad 2(k-1)\le q.
\end{equation}
For $1\le j<k$, there are exactly $k-j$ index pairs separated by $j$.
Under~\eqref{eq:half}, their distance classes $|jg|_n$ are distinct:
a coincidence would give $j\equiv\pm\ell\pmod q$; for distinct
$j,\ell\in\{1,\ldots,k-1\}$, neither $j-\ell$ nor $j+\ell$ can be a
multiple of $q$. The points are distinct as well. Thus the constituent
has multiplicities $1,\ldots,k-1$, including at the endpoint in~\eqref{eq:half}.

\newpage
\section{A counterexample in $\Z_{34}$}
\begin{theorem}\label{thm:34}
The family
\[
 \mathcal F_{34}=\{\AP_{34}(1,10),\AP_{34}(6,9),\AP_{34}(12,5)\}
\]
is primitive and Erd\H{o}s-deep with parameter $K=14$.
All its constituents are individually Erd\H{o}s-deep.
\end{theorem}
\begin{proof}
The actual point sets, written in progression order, are
\begin{align*}
 A&=\{0,1,2,3,4,5,6,7,8,9\},\\
 B&=\{0,6,12,18,24,30,2,8,14\},\\
 C&=\{0,12,24,2,14\}.
\end{align*}
Their generator orders are $34,17,17$. Condition~\eqref{eq:half}
reads $18\le34$, $16\le17$, and $8\le17$, so each constituent has
its asserted cardinality and is individually Erd\H{o}s-deep.
The first generator is $1$, giving primitivity.

For $B$, the distances for index gaps $1,\ldots,8$ and their weights are
\[
 (6,12,16,10,4,2,8,14),\qquad (8,7,6,5,4,3,2,1),
\]
respectively.
For $C$, the distances are $(12,10,2,14)$, with weights $(4,3,2,1)$.
Combining these with the interval $A$ gives every entry below.
\begin{center}
\renewcommand{\arraystretch}{1.08}
\begin{tabular}{rrrrr}
\toprule
Distance $d$ & $m_A(d)$ & $m_B(d)$ & $m_C(d)$ & $m_{\mathcal F_{34}}(d)$\\
\midrule
1&9&0&0&9\\
2&8&3&2&13\\
3&7&0&0&7\\
4&6&4&0&10\\
5&5&0&0&5\\
6&4&8&0&12\\
7&3&0&0&3\\
8&2&2&0&4\\
9&1&0&0&1\\
10&0&5&3&8\\
12&0&7&4&11\\
14&0&1&1&2\\
16&0&6&0&6\\
\bottomrule
\end{tabular}
\end{center}
All omitted distances have multiplicity zero. The last column is a
permutation of $1,\ldots,13$. As a separate total-count check,
$\binom{10}{2}+\binom92+\binom52=45+36+10=91=\binom{14}{2}$.
\end{proof}

\newpage
\section{A counterexample in $\Z_{48}$}
\begin{theorem}\label{thm:48}
The family
\[
 \mathcal F_{48}=\{\AP_{48}(1,20),\AP_{48}(6,5),\AP_{48}(18,5)\}
\]
is primitive and Erd\H{o}s-deep with parameter $K=21$.
All its constituents are individually Erd\H{o}s-deep.
\end{theorem}
\begin{proof}
The constituents are
\[
 A=\{0,1,\ldots,19\},\quad
 B=\{0,6,12,18,24\},\quad C=\{0,18,36,6,24\}.
\]
Their generator orders are $48,8,8$; the three instances
of~\eqref{eq:half} are $38\le48$, $8\le8$, and $8\le8$.
They are individually Erd\H{o}s-deep, and the first generator again
ensures primitivity.

The interval $A$ gives $m_A(d)=20-d$ for $1\le d\le19$, and zero
otherwise. For $B$, the four index-gap distances are $(6,12,18,24)$;
for $C$, they are $(18,12,6,24)$. Both have index-gap weights
$(4,3,2,1)$, so their combined contributions are
\begin{center}
\begin{tabular}{rrrrr}
\toprule
$d$ & 6&12&18&24\\
\midrule
$m_B(d)$&4&3&2&1\\
$m_C(d)$&2&3&4&1\\
$m_B(d)+m_C(d)$&6&6&6&2\\
\bottomrule
\end{tabular}
\end{center}
Consequently the multiplicities at distances $6,12,18$ change from
$14,8,2$ to $20,14,8$, respectively, and distance $24$ supplies the
new multiplicity $2$. All other multiplicities of $A$ stay unchanged.
The result is exactly $1,\ldots,20$, each once. The pair-count check is
$\binom{20}{2}+2\binom52=190+20=210=\binom{21}{2}$.
\end{proof}

\section{Consequence and exact source scope}
The fifteen triples in the printed classification~\cite[Conjecture~1,
p.~14]{DG} do not include $(10,9,5)$ or $(20,5,5)$. Theorems
\ref{thm:34} and~\ref{thm:48} therefore refute its necessity claim,
even after imposing primitivity, minimum constituent length three,
and individual Erd\H{o}s-deepness.

The endpoint issue must be separated from the exact computations.
The paper's Theorem~1~\cite[p.~3]{DG} includes the $+1$ endpoint
in its characterization of individually Erd\H{o}s-deep APs. Gaede's standing convention~\cite[\S2.2, p.~16]{G} allows
$k_i\le\lfloor\ord(g_i)/2\rfloor+1$. Conjecture~4.5
\cite[pp.~14,42]{G} also displays the $+1$ bound for the largest
constituent. An earlier general convention~\cite[Remark~1.10,
p.~5]{G} instead omits $+1$; applying that stricter convention to
every constituent excludes both examples. The displayed search
pseudocode~\cite[p.~43]{G} explicitly bounds only $k_1$ using $g_1$.
Both examples satisfy that displayed bound. Thus the pseudocode
\emph{does not establish why the examples were omitted}, and the
implementation cause remains unknown. The all-constituent
strict-half-order variant is neither proved nor disproved here.

\newpage
\section{Formal verification and reproducibility}
The accompanying \texttt{ErdosDeepCounterexamples.lean} defines
modular point lists, the cyclic metric, and a structurally recursive
list of unordered-pair distances. On an input $x::L$, that list
contains the distances from $x$ to every element of $L$, followed
by the pair-distance list of $L$. Distinctness of every example's
point list is separately checked, so positions correspond to
actual distinct elements.

For each $d=1,\ldots,\lfloor n/2\rfloor$, the development counts
occurrences of $d$. It removes zero counts and sorts the remaining
multiplicities by insertion sort. Equality with the list
$[1,\ldots,K-1]$ is the exact depth predicate. The following are
kernel-checked:
\begin{itemize}
\item primitive gcd, ordering, minimum lengths, and the
half-order bounds;
\item distinct point lists and the complete histograms;
\item individual depth of all six constituents and total depth of
both families;
\item exclusion from the fifteen listed triples and the negation
of an explicitly normalized universal classification statement.
\end{itemize}
The checks use Lean~4.32.2, mathlib commit
\begin{center}\small\texttt{905b95818eb32af7874a58b427f50c1711a5e96c}.\end{center}
There is no \texttt{sorry}, \texttt{native\_decide}, or added axiom.
The histogram and individual-depth theorems have no axioms;
the bundled counterexamples and negated classification use only
the standard \texttt{propext} and \texttt{Quot.sound} principles.
The script \texttt{check-lean.sh} rebuilds the certificate, and
\texttt{lean-build.log} records the axiom audit.

A C++ search through $n\le60$ produced the fifteen listed length
triples and these two additional triples. A separate Python
program enumerated unordered pairs of the actual modular points
and checked all 649 recorded representative hits. This bounded
search is supplemental evidence, not a classification theorem;
neither counterexample depends on its exhaustiveness.

\paragraph{Prior-art check.}
A limited public-web search on 14 October 2026 used the paper's
exact title, the authors and triple-classification terminology,
and the two extra length triples in spaced and unspaced forms.
It located the original article, thesis, and authors' research
pages, but no relevant earlier report of these examples. That
search is not exhaustive and does not establish priority. Its
queries and checked primary URLs are recorded in the package.

\begin{thebibliography}{9}
\bibitem{DG} P.~J.~Dukes and T.~Gaede,
\emph{Families of modular arithmetic progressions with an interval
of distance multiplicities}, Integers \textbf{23} (2023), \#A25.
\url{https://math.colgate.edu/~integers/x25/x25.pdf}.
\bibitem{G} T.~Gaede,
\emph{Erd\H{o}s-Deep Families of Arithmetic Progressions},
M.Sc. thesis, University of Victoria, 2022.
\url{https://hdl.handle.net/1828/14160}.
\end{thebibliography}
\end{document}
